% =============================================================================
% Aula 05 — Relacionamentos 1:N e N:N · Categoria no laboratorio · Cliente e venda no StockSales
% Beamer + tema Madrid | Curso Entra21 — Spring Framework
% Compilar: tectonic -X compile slides-aula-05.tex
% =============================================================================
\documentclass[aspectratio=169,11pt]{beamer}

\usetheme{Madrid}
\usecolortheme{default}

\usepackage[T1]{fontenc}
\usepackage[utf8]{inputenc}
\usepackage[brazil]{babel}
\usepackage{booktabs}
\usepackage{graphicx}
\usepackage{hyperref}
\usepackage{listings}
\usepackage{xcolor}
\usepackage{tikz}
\usetikzlibrary{arrows.meta, positioning, shapes.geometric}

\definecolor{codebg}{RGB}{245,245,245}
\definecolor{codegreen}{RGB}{40,120,40}
\definecolor{codeblue}{RGB}{30,70,160}
\definecolor{codered}{RGB}{180,50,50}

\lstset{
  basicstyle=\ttfamily\small,
  backgroundcolor=\color{codebg},
  frame=single,
  rulecolor=\color{black!20},
  breaklines=true,
  columns=fullflexible,
  keepspaces=true,
  showstringspaces=false,
  keywordstyle=\color{codeblue}\bfseries,
  commentstyle=\color{codegreen},
  stringstyle=\color{codered},
  tabsize=2,
  literate={á}{{\'a}}1 {é}{{\'e}}1 {í}{{\'i}}1 {ó}{{\'o}}1 {ú}{{\'u}}1
           {Á}{{\'A}}1 {É}{{\'E}}1 {Í}{{\'I}}1 {Ó}{{\'O}}1 {Ú}{{\'U}}1
           {ã}{{\~a}}1 {õ}{{\~o}}1 {Ã}{{\~A}}1 {Õ}{{\~O}}1
           {ç}{{\c{c}}}1 {Ç}{{\c{C}}}1
           {ê}{{\^e}}1 {Ê}{{\^E}}1 {â}{{\^a}}1 {Â}{{\^A}}1
           {à}{{\`a}}1 {À}{{\`A}}1
}

\newcommand{\onde}[1]{%
  \begin{block}{Onde estamos}
    #1
  \end{block}
}

\title{Aula 05 --- Spring Framework}
\subtitle{Relacionamentos · 1{:}N · N{:}N · Cliente e venda}
\author{Mauricio Duailibi Neto}
\institute{Turma Java Entra21}
\date{}

\begin{document}

\begin{frame}
  \titlepage
\end{frame}

\begin{frame}{Como será a aula de hoje}
  \begin{block}{Parte 1 --- 18h30 às 20h}
    \begin{itemize}
      \item Revisão da Aula 04
      \item Relacionamento: um cadastro aponta para outro
      \item 1{:}N com categoria e material
      \item N{:}N: a ideia da lista do meio
      \item Projeto \textbf{laboratorio} --- fazemos juntos
    \end{itemize}
  \end{block}

  \vspace{0.4em}
  \begin{block}{Parte 2 --- 20h20 às 22h}
    \begin{itemize}
      \item Voltamos ao \textbf{StockSales}
      \item Segundo cadastro: cliente
      \item Primeira venda: um cliente e um produto
    \end{itemize}
  \end{block}
\end{frame}

\begin{frame}{Dois projetos nesta aula}
  Não vamos misturar os códigos.

  \vspace{0.6em}
  \begin{description}
    \item[laboratorio] Primeira parte da aula.
      Aqui aprendemos a ligar categoria e material.
    \item[StockSales] O sistema da loja, que já existe.
      Só depois do intervalo.
  \end{description}

  \vspace{0.6em}
  Quando aparecer um código, o slide vai dizer:\\
  \textbf{qual projeto} e \textbf{qual arquivo} estamos mexendo.
\end{frame}

\begin{frame}{Mapa da aula}
  \tableofcontents
\end{frame}

% #############################################################################
\section{Revisão da Aula 04}

\begin{frame}{O que já sabemos}
  Nos dois projetos, até agora, existe \textbf{um} model
  e \textbf{um} Controller.

  \vspace{0.5em}
  \begin{itemize}
    \item laboratorio: \texttt{Material} --- GET, POST, PUT, DELETE
    \item StockSales: \texttt{Produto} --- GET, POST, PUT, DELETE
    \item A lista fica \textbf{na classe} do Controller
    \item PUT altera pelo id. Não gera id novo. Não dá \texttt{add}
    \item O JavaScript chama a API com Axios
  \end{itemize}

  \vspace{0.5em}
  Cada cadastro vive sozinho.\\
  O material não aponta para ninguém. O produto tampouco.
\end{frame}

\begin{frame}{Os quatro verbos}
  \begin{center}
  \begin{tabular}{lll}
    \toprule
    Ação & Verbo HTTP & O que faz \\
    \midrule
    Listar / ler & GET & Lê, não muda nada \\
    Cadastrar & POST & Cria um item novo \\
    Editar & PUT & Troca os dados de um item \\
    Excluir & DELETE & Apaga um item \\
    \bottomrule
  \end{tabular}
  \end{center}

  \vspace{0.7em}
  Hoje esses verbos continuam iguais.\\
  O que muda é o \textbf{conteúdo} do cadastro:\\
  um objeto passa a guardar o id de outro.
\end{frame}

\begin{frame}{O que ainda falta}
  No laboratorio, o material tem id, nome e quantidade.\\
  Ele não diz de que grupo é.

  \vspace{0.5em}
  No StockSales, só existe produto.\\
  A loja ainda não tem cliente e ainda não registra venda.

  \vspace{0.6em}
  \begin{block}{Pergunta de hoje}
    Como um cadastro aponta para outro\\
    sem copiar todos os dados?
  \end{block}
\end{frame}

% #############################################################################
\section{Relacionamentos}

\begin{frame}{Cadastros que se falam}
  Até agora cada lista vivia sozinha:\\
  materiais de um lado, produtos de outro.

  \vspace{0.5em}
  No mundo real eles se encaixam:
  \begin{itemize}
    \item O material \textbf{pertence} a uma categoria
    \item Depois, na loja, a venda \textbf{pertence} a um cliente
    \item E aponta para um produto
  \end{itemize}

  \vspace{0.5em}
  Essa ligação chama \textbf{relacionamento}.
\end{frame}

\begin{frame}{Uma analogia que todo mundo já viu}
  Uma turma tem vários alunos.\\
  Um aluno está em \textbf{uma} turma.

  \vspace{0.5em}
  Não copiamos a lista de alunos dentro da turma\\
  toda vez que alguém pergunta o nome da turma.

  \vspace{0.5em}
  Guardamos no aluno o número da turma:\\
  \texttt{turmaId = 2}.

  \vspace{0.5em}
  Quando precisamos do nome, procuramos a turma 2\\
  na lista de turmas.
\end{frame}

\begin{frame}{Dois jeitos de ligar}
  \begin{center}
  \begin{tabular}{ll}
    \toprule
    Nome & Significado \\
    \midrule
    1:N & Um de um lado, vários do outro \\
    N:N & Vários dos dois lados \\
    \bottomrule
  \end{tabular}
  \end{center}

  \vspace{0.7em}
  1:N --- uma categoria, vários materiais.\\
  Cada material tem \textbf{uma} categoria.

  \vspace{0.5em}
  N:N --- um material tem várias etiquetas,\\
  e a mesma etiqueta serve para vários materiais.
\end{frame}

\begin{frame}{1{:}N --- o desenho}
  Uma categoria aponta para vários materiais.\\
  O material aponta para \textbf{uma} categoria.

  \vspace{0.6em}
  \begin{center}
  \begin{tikzpicture}[
      box/.style={draw, rounded corners, align=center,
                  minimum width=2.6cm, minimum height=0.85cm,
                  font=\small},
      seta/.style={-{Latex}, thick}
    ]
    \node[box, fill=blue!15] (cat) {Papelaria};
    \node[box, fill=green!20, right=2.2cm of cat] (c1) {Caneta};
    \node[box, fill=green!20, below=0.35cm of c1] (c2) {Caderno};
    \node[box, fill=green!20, below=0.35cm of c2] (c3) {Grampeador};

    \draw[seta] (cat.east) -- ++(0.5,0) |- (c1.west);
    \draw[seta] (cat.east) -- ++(0.5,0) |- (c2.west);
    \draw[seta] (cat.east) -- ++(0.5,0) |- (c3.west);
  \end{tikzpicture}
  \end{center}

  \vspace{0.3em}
  Papelaria é o ``1''.\\
  Caneta, Caderno e Grampeador são o ``N''.
\end{frame}

\begin{frame}{Como gravamos o 1{:}N}
  Não colocamos uma lista de materiais dentro da categoria.

  \vspace{0.5em}
  Colocamos no material o id da categoria:

  \vspace{0.4em}
  \begin{center}
  \begin{tabular}{lll}
    \toprule
    Material & quantidade & categoriaId \\
    \midrule
    Caneta & 10 & 1 \\
    Caderno & 5 & 1 \\
    Mouse & 3 & 2 \\
    \bottomrule
  \end{tabular}
  \end{center}

  \vspace{0.5em}
  A categoria 1 é Papelaria.\\
  A categoria 2 é Informática.

  \vspace{0.4em}
  É o mesmo jeito do \texttt{turmaId} no aluno.
\end{frame}

\begin{frame}[fragile]{O JSON do material com categoria}
  O cadastro passa a levar mais um número:

  \vspace{0.4em}
\begin{lstlisting}[basicstyle=\ttfamily\footnotesize]
{"nome": "Caneta", "quantidade": 10, "categoriaId": 1}
\end{lstlisting}

  \vspace{0.5em}
  O servidor não precisa do texto ``Papelaria'' nesse JSON.\\
  O id 1 já aponta para ela.

  \vspace{0.5em}
  Se a pessoa mudar o nome da categoria depois,\\
  a Caneta continua ligada ao id 1.\\
  O nome novo aparece sozinho.
\end{frame}

\begin{frame}{N{:}N --- por que um id só não basta}
  Um material pode ser Frágil \textbf{e} Uso interno.\\
  A etiqueta Uso interno serve para Caneta \textbf{e} Caderno.

  \vspace{0.5em}
  Se o material tiver um único \texttt{etiquetaId},\\
  ele só cabe em uma etiqueta.

  \vspace{0.5em}
  Se a etiqueta tiver um único \texttt{materialId},\\
  ela só cabe em um material.

  \vspace{0.5em}
  Os dois lados são ``vários''.\\
  Um campo só não resolve.
\end{frame}

\begin{frame}{N{:}N --- a lista do meio}
  Criamos uma terceira lista.\\
  Cada linha é um par: este material + esta etiqueta.

  \vspace{0.4em}
  \begin{center}
  \begin{tabular}{ll}
    \toprule
    materialId & etiquetaId \\
    \midrule
    1 (Caneta) & 1 (Fragil) \\
    1 (Caneta) & 2 (Uso interno) \\
    2 (Caderno) & 2 (Uso interno) \\
    \bottomrule
  \end{tabular}
  \end{center}

  \vspace{0.5em}
  Caneta aparece duas vezes: tem duas etiquetas.\\
  Uso interno aparece duas vezes: serve para dois materiais.
\end{frame}

\begin{frame}{N{:}N --- o desenho}
  \begin{center}
  \begin{tikzpicture}[
      box/.style={draw, rounded corners, align=center,
                  minimum width=2.4cm, minimum height=0.8cm,
                  font=\small},
      meio/.style={draw, rounded corners, align=center,
                   minimum width=1.6cm, minimum height=0.55cm,
                   font=\scriptsize},
      seta/.style={-{Latex}, thick}
    ]
    \node[box, fill=green!20] (m1) {Caneta};
    \node[box, fill=green!20, below=0.7cm of m1] (m2) {Caderno};

    \node[meio, fill=yellow!25, right=1.5cm of m1] (l1) {1 -- 1};
    \node[meio, fill=yellow!25, below=0.15cm of l1] (l2) {1 -- 2};
    \node[meio, fill=yellow!25, below=0.15cm of l2] (l3) {2 -- 2};

    \node[box, fill=blue!15, right=1.5cm of l1] (e1) {Fragil};
    \node[box, fill=blue!15, below=0.7cm of e1] (e2) {Uso interno};

    \draw[seta] (m1) -- (l1);
    \draw[seta] (m1) -- (l2);
    \draw[seta] (m2) -- (l3);
    \draw[seta] (l1) -- (e1);
    \draw[seta] (l2) -- (e2);
    \draw[seta] (l3) -- (e2);
  \end{tikzpicture}
  \end{center}

  \vspace{0.3em}
  O amarelo é a lista do meio.\\
  Sem ela, os dois lados não se encontram.
\end{frame}

\begin{frame}{Depois do intervalo, a mesma ideia na loja}
  Até hoje o StockSales só tem \texttt{Produto}.

  \vspace{0.5em}
  Na segunda parte a loja ganha \textbf{cliente} e \textbf{venda}.\\
  A venda guarda \texttt{clienteId} e \texttt{produtoId}:\\
  o mesmo 1{:}N da categoria com o material.

  \vspace{0.5em}
  N{:}N (vários produtos na mesma venda) fica como extra,\\
  se der tempo.
\end{frame}

\begin{frame}{O que não vamos usar hoje}
  \begin{itemize}
    \item Banco de dados
    \item Anotações do tipo \texttt{@OneToMany}
    \item JPA
  \end{itemize}

  \vspace{0.5em}
  Continuamos com lista na memória e um \texttt{int}\\
  guardando o id do outro cadastro.

  \vspace{0.5em}
  O raciocínio é o mesmo que o banco vai usar depois.\\
  Só a ferramenta muda.
\end{frame}

% #############################################################################
\section{Projeto laboratorio}

\begin{frame}{O projeto laboratorio}
  Continuamos no projeto \textbf{laboratorio}.\\
  Não geramos outro no \texttt{start.spring.io}.

  \vspace{0.5em}
  \begin{block}{O que já deve estar pronto}
    Material com GET, POST, PUT e DELETE.\\
    Formulário na tela. Botões Editar e Excluir.
  \end{block}

  \vspace{0.5em}
  Hoje entra um cadastro novo: \textbf{categoria}.\\
  O material ganha um \texttt{categoriaId}.
\end{frame}

\begin{frame}{Passo 1 --- Subir o Spring}
  \onde{Projeto \textbf{laboratorio}}

  No terminal, dentro da pasta \texttt{laboratorio}:

  \vspace{0.4em}
  \texttt{./mvnw spring-boot:run}\\
  {\small No Windows: \texttt{.\textbackslash mvnw.cmd spring-boot:run}}

  \vspace{0.5em}
  Abram no navegador: \texttt{http://localhost:8080}

  \vspace{0.5em}
  \begin{alertblock}{Um Spring por vez}
    Se o StockSales ainda estiver na porta 8080,
    parem o outro antes.
  \end{alertblock}
\end{frame}

\begin{frame}{Quais arquivos já existem}
  Projeto: \textbf{laboratorio}

  \vspace{0.4em}
  \begin{center}
  \begin{tabular}{ll}
    \toprule
    Arquivo & Para quê \\
    \midrule
    \texttt{Material.java} & Os dados de um material \\
    \texttt{MaterialController.java} & GET, POST, PUT, DELETE \\
    \texttt{static/index.html} & A página \\
    \texttt{static/js/app.js} & O JavaScript da página \\
    \bottomrule
  \end{tabular}
  \end{center}

  \vspace{0.5em}
  Vamos \textbf{acrescentar} categoria e a classe \texttt{Dados}.
\end{frame}

\begin{frame}[fragile]{Onde cada arquivo fica}
\begin{lstlisting}[basicstyle=\ttfamily\footnotesize]
laboratorio/
  src/main/java/com/entra21/laboratorio/
    LaboratorioApplication.java
    Dados.java                  (novo)
    model/Material.java
    model/Categoria.java        (novo)
    controller/MaterialController.java
    controller/CategoriaController.java  (novo)
  src/main/resources/static/
    index.html
    js/app.js
\end{lstlisting}
\end{frame}

\begin{frame}{Por que a classe Dados}
  Vamos ter dois Controllers: material e categoria.

  \vspace{0.5em}
  Se cada um guardar a própria lista, um \textbf{não vê}\\
  o cadastro do outro.

  \vspace{0.5em}
  Por isso uma classe só, \texttt{Dados}, com as listas
  \texttt{static}.

  \vspace{0.5em}
  \texttt{static} = uma lista só, visível para qualquer\\
  classe do projeto.\\
  Até agora a lista ficava dentro do Controller.\\
  Hoje ela sai de lá, porque passou a ter dois Controllers.
\end{frame}

\begin{frame}[fragile]{Classe Dados}
  \onde{Projeto \textbf{laboratorio}\\
        Arquivo novo:\\
        \texttt{src/main/java/com/entra21/laboratorio/Dados.java}}

\begin{lstlisting}[language=Java,basicstyle=\ttfamily\scriptsize]
package com.entra21.laboratorio;

import java.util.ArrayList;
import com.entra21.laboratorio.model.Material;
import com.entra21.laboratorio.model.Categoria;

public class Dados {
  public static ArrayList<Material> materiais = new ArrayList<>();
  public static int proximoIdMaterial = 1;
  public static ArrayList<Categoria> categorias = new ArrayList<>();
  public static int proximoIdCategoria = 1;
}
\end{lstlisting}
\end{frame}

\begin{frame}[fragile]{Mover a lista de materiais}
  \onde{\texttt{MaterialController.java}}

  Tirem os campos \texttt{materiais} e \texttt{proximoId}\\
  de dentro da classe.

  \vspace{0.4em}
  Em todos os métodos, troquem:
  \begin{itemize}
    \item \texttt{materiais} $\rightarrow$ \texttt{Dados.materiais}
    \item \texttt{proximoId} $\rightarrow$ \texttt{Dados.proximoIdMaterial}
  \end{itemize}

  \vspace{0.4em}
  O import:
  \texttt{com.entra21.laboratorio.Dados}

  \vspace{0.4em}
  Confiram GET, POST, PUT e DELETE no navegador\\
  \textbf{antes} de seguir. A lista precisa continuar funcionando.
\end{frame}

% #############################################################################
\section{Categoria (1{:}N)}

\begin{frame}[fragile]{Classe Categoria}
  \onde{Arquivo novo:\\
        \texttt{src/main/java/com/entra21/laboratorio/model/Categoria.java}}

\begin{lstlisting}[language=Java,basicstyle=\ttfamily\footnotesize]
public class Categoria {
  private int id;
  private String nome;

  public Categoria() {
  }

  public Categoria(int id, String nome) {
    this.id = id;
    this.nome = nome;
  }
}
\end{lstlisting}

  Completem os gets e sets de \texttt{id} e \texttt{nome}.
\end{frame}

\begin{frame}[fragile]{categoriaId no Material}
  \onde{\texttt{model/Material.java}}

  Acrescentem o campo e os get/set.

\begin{lstlisting}[language=Java,basicstyle=\ttfamily\footnotesize]
private int categoriaId;

public int getCategoriaId() {
  return categoriaId;
}

public void setCategoriaId(int categoriaId) {
  this.categoriaId = categoriaId;
}
\end{lstlisting}

  O material agora aponta para uma categoria.
\end{frame}

\begin{frame}[fragile]{Construtor cheio do Material}
  \onde{Ainda em \texttt{Material.java}}

  Se ainda não tiverem, acrescentem este construtor.\\
  Mantenham o vazio. O vazio continua obrigatório no POST.

\begin{lstlisting}[language=Java,basicstyle=\ttfamily\footnotesize]
public Material(int id, String nome,
                int quantidade, int categoriaId) {
  this.id = id;
  this.nome = nome;
  this.quantidade = quantidade;
  this.categoriaId = categoriaId;
}
\end{lstlisting}
\end{frame}

\begin{frame}[fragile,shrink=5]{CategoriaController --- seed e GET}
  \onde{\texttt{controller/CategoriaController.java}, arquivo novo}

\begin{lstlisting}[language=Java,basicstyle=\ttfamily\scriptsize]
@RestController
public class CategoriaController {

  public CategoriaController() {
    if (Dados.categorias.isEmpty()) {
      Dados.categorias.add(new Categoria(1, "Papelaria"));
      Dados.categorias.add(new Categoria(2, "Informatica"));
      Dados.categorias.add(new Categoria(3, "Limpeza"));
      Dados.proximoIdCategoria = 4;
    }
  }

  @GetMapping("/api/categorias")
  public ArrayList<Categoria> listar() {
    return Dados.categorias;
  }
}
\end{lstlisting}
\end{frame}

\begin{frame}[fragile]{POST de categoria}
  \onde{Ainda em \texttt{CategoriaController.java}}

\begin{lstlisting}[language=Java,basicstyle=\ttfamily\footnotesize]
@PostMapping("/api/categorias")
public Categoria cadastrar(@RequestBody Categoria categoria) {
  categoria.setId(Dados.proximoIdCategoria);
  Dados.proximoIdCategoria = Dados.proximoIdCategoria + 1;
  Dados.categorias.add(categoria);
  return categoria;
}
\end{lstlisting}

  \vspace{0.4em}
  Testem no navegador:\\
  \texttt{http://localhost:8080/api/categorias}

  \vspace{0.3em}
  Devem aparecer Papelaria, Informatica e Limpeza.
\end{frame}

\begin{frame}[fragile]{PUT de material também grava a categoria}
  \onde{\texttt{MaterialController.java}, método \texttt{atualizar}}

  Junto de \texttt{setNome} e \texttt{setQuantidade}:

\begin{lstlisting}[language=Java,basicstyle=\ttfamily\footnotesize]
Dados.materiais.get(i).setNome(dados.getNome());
Dados.materiais.get(i).setQuantidade(dados.getQuantidade());
Dados.materiais.get(i).setCategoriaId(dados.getCategoriaId());
\end{lstlisting}

  \vspace{0.4em}
  Sem essa linha, o Editar apaga a categoria\\
  e deixa \texttt{categoriaId} em 0.
\end{frame}

\begin{frame}{O POST de material já lê o campo novo}
  Não precisamos de método extra.

  \vspace{0.5em}
  O JSON agora traz \texttt{categoriaId}.\\
  O Spring preenche o campo no objeto.\\
  O \texttt{add} na lista já guarda tudo.

  \vspace{0.5em}
  Só o JavaScript precisa passar a \textbf{enviar} esse número.
\end{frame}

\begin{frame}{Conferindo o backend}
  \begin{itemize}
    \item [ ] \texttt{/api/categorias} lista as três categorias
    \item [ ] \texttt{/api/materiais} continua listando
    \item [ ] Um POST de material com \texttt{categoriaId} no JSON
          devolve o campo preenchido
  \end{itemize}

  \vspace{0.5em}
  Se \texttt{/api/categorias} der 404, o Controller novo\\
  não compilou ou o Spring não reiniciou.
\end{frame}

% #############################################################################
\section{Tela com select}

\begin{frame}{A pessoa não precisa decorar o id}
  Digitar \texttt{categoriaId = 2} na mão é fácil de errar.

  \vspace{0.5em}
  No HTML existe o \texttt{<select>}: uma lista para escolher.

  \vspace{0.5em}
  A pessoa vê o \textbf{nome} (Papelaria).\\
  O JavaScript manda o \textbf{id} (1) no JSON.
\end{frame}

\begin{frame}[fragile]{O select no HTML}
  \onde{Projeto \textbf{laboratorio}\\
        \texttt{static/index.html}, dentro do form}

\begin{lstlisting}[basicstyle=\ttfamily\footnotesize]
<label>Categoria</label>
<select id="categoriaId"></select>
\end{lstlisting}

  \vspace{0.5em}
  Fica junto dos inputs de nome e quantidade.\\
  O JS vai encher as opções depois do GET.
\end{frame}

\begin{frame}[fragile]{Uma lista de categorias no JS}
  \onde{\texttt{static/js/app.js}, no topo do arquivo}

  Vamos guardar a resposta do GET.\\
  A lista da tela usa essa variável para achar o nome.

\begin{lstlisting}[basicstyle=\ttfamily\footnotesize]
var categorias = [];
\end{lstlisting}
\end{frame}

\begin{frame}[fragile]{Preencher o select}
  \onde{\texttt{static/js/app.js}}

\begin{lstlisting}[basicstyle=\ttfamily\scriptsize]
function carregarCategorias() {
  axios.get("/api/categorias").then(function (resposta) {
    categorias = resposta.data;
    var select = document.getElementById("categoriaId");
    select.innerHTML = "";
    for (var i = 0; i < categorias.length; i++) {
      var opcao = document.createElement("option");
      opcao.value = categorias[i].id;
      opcao.textContent = categorias[i].nome;
      select.appendChild(opcao);
    }
  });
}
\end{lstlisting}

  \texttt{value} é o id. O texto visível é o nome.
\end{frame}

\begin{frame}{Por que \texttt{value} é o id}
  O que a pessoa vê: Papelaria.

  \vspace{0.5em}
  O que o servidor precisa: \texttt{1}.

  \vspace{0.5em}
  O \texttt{<option>} guarda os dois:\\
  \texttt{value="1"} e o texto Papelaria.

  \vspace{0.5em}
  \texttt{select.value} devolve \texttt{"1"} --- ainda é texto.\\
  Por isso usamos \texttt{Number(...)} na hora de montar o JSON.
\end{frame}

\begin{frame}[fragile]{Achar o nome pelo id}
  \onde{\texttt{static/js/app.js}}

\begin{lstlisting}[basicstyle=\ttfamily\footnotesize]
function nomeDaCategoria(categoriaId) {
  for (var i = 0; i < categorias.length; i++) {
    if (categorias[i].id === categoriaId) {
      return categorias[i].nome;
    }
  }
  return "";
}
\end{lstlisting}

  Percorre a lista. Quando o id bate, devolve o nome.\\
  Se não achar, devolve texto vazio.
\end{frame}

\begin{frame}[fragile]{Mostrar o nome na lista}
  \onde{\texttt{static/js/app.js}, no \texttt{for} de
        \texttt{carregarMateriais}}

  Troquem o \texttt{textContent} do \texttt{li} por:

\begin{lstlisting}[basicstyle=\ttfamily\footnotesize]
li.textContent = material.nome
  + " - quantidade: " + material.quantidade
  + " - " + nomeDaCategoria(material.categoriaId)
  + " ";
\end{lstlisting}

  A pessoa lê ``Caneta - quantidade: 10 - Papelaria''.\\
  Não lê o número 1.
\end{frame}

\begin{frame}[fragile]{O submit manda o categoriaId}
  \onde{\texttt{static/js/app.js}, no objeto \texttt{dados}\\
        do submit --- no POST e no PUT}

\begin{lstlisting}[basicstyle=\ttfamily\footnotesize]
var dados = {
  nome: document.getElementById("nome").value,
  quantidade: Number(
    document.getElementById("quantidade").value),
  categoriaId: Number(
    document.getElementById("categoriaId").value)
};
\end{lstlisting}

  Sem \texttt{Number}, o Java pode recusar o campo,\\
  porque espera \texttt{int}.
\end{frame}

\begin{frame}[fragile]{Editar também preenche o select}
  \onde{\texttt{static/js/app.js}, no clique de Editar}

  Além de nome e quantidade:

\begin{lstlisting}[basicstyle=\ttfamily\footnotesize]
document.getElementById("categoriaId").value =
  evento.target.getAttribute("data-categoria-id");
\end{lstlisting}

  \vspace{0.4em}
  No botão Editar, acrescentem:

\begin{lstlisting}[basicstyle=\ttfamily\footnotesize]
btnEditar.setAttribute("data-categoria-id",
  material.categoriaId);
\end{lstlisting}
\end{frame}

\begin{frame}[fragile]{Quando carregar a página}
  \onde{No final de \texttt{app.js}}

  As categorias precisam chegar \textbf{antes} da lista,\\
  senão \texttt{nomeDaCategoria} não acha o nome.

\begin{lstlisting}[basicstyle=\ttfamily\footnotesize]
carregarCategorias();
carregarMateriais();
\end{lstlisting}

  \vspace{0.4em}
  Se o nome da categoria sair vazio na primeira carga,\\
  chamem \texttt{carregarMateriais} \textbf{dentro} do
  \texttt{then} de \texttt{carregarCategorias}.
\end{frame}

\begin{frame}{O caminho do cadastro com categoria}
  \begin{enumerate}
    \item A página pede GET \texttt{/api/categorias}
    \item O select mostra Papelaria, Informatica, Limpeza
    \item A pessoa escolhe Papelaria, digita Caneta, envia
    \item O JS manda \texttt{categoriaId: 1} no POST
    \item O material entra na lista com esse número
    \item Na tela, o JS troca 1 por Papelaria
  \end{enumerate}

  \vspace{0.4em}
  No Network: um GET de categorias e um POST de material.
\end{frame}

\begin{frame}{Conferindo o 1{:}N}
  \begin{itemize}
    \item [ ] O select mostra as três categorias
    \item [ ] Cadastrar Caneta em Papelaria mostra o nome
          na lista, não só o id
    \item [ ] Editar troca a categoria e não duplica o material
    \item [ ] Mouse em Informatica não aparece como Papelaria
    \item [ ] No Network, o JSON do POST tem \texttt{categoriaId}
  \end{itemize}
\end{frame}

\begin{frame}{Erros que mais aparecem}
  \begin{itemize}
    \item Select vazio --- \texttt{carregarCategorias} não rodou
    \item Nome da categoria em branco --- a lista de materiais
          rodou antes do GET de categorias
    \item \texttt{categoriaId} some no Editar --- faltou o
          \texttt{setCategoriaId} no PUT
    \item Lista de materiais some --- ainda está no Controller,
          não em \texttt{Dados}
    \item Porta 8080 ocupada
  \end{itemize}
\end{frame}

% #############################################################################
\section{Desafios}

\begin{frame}{Agora é com vocês}
  Continuem no projeto \textbf{laboratorio}.\\
  Mesmos arquivos.

  \vspace{0.5em}
  \begin{enumerate}
    \item Filtrar a lista por categoria
    \item Não apagar categoria que ainda tem material
    \item Etiquetas: o N:N
  \end{enumerate}

  \vspace{0.5em}
  Salve $\rightarrow$ reinicie se precisar $\rightarrow$ teste no navegador.\\
  F12, aba Network.
\end{frame}

\begin{frame}{Desafio 1 --- filtro por categoria}
  \onde{\texttt{index.html} e \texttt{static/js/app.js}}

  Coloquem um segundo select, por exemplo\\
  \texttt{id="filtro-categoria"}.

  \vspace{0.4em}
  A primeira opção: \texttt{value=""} e texto ``Todas''.\\
  As outras: as mesmas categorias.

  \vspace{0.4em}
  Em \texttt{carregarMateriais}, se o filtro tiver um id,\\
  só mostrem o material cujo \texttt{categoriaId} é aquele.

  \vspace{0.4em}
  No \texttt{change} do select, chamem
  \texttt{carregarMateriais} de novo.
\end{frame}

\begin{frame}[fragile]{Desafio 1 --- o if do filtro}
  Ideia do \texttt{for} (vocês escrevem o resto):

\begin{lstlisting}[basicstyle=\ttfamily\footnotesize]
var filtro = document.getElementById(
  "filtro-categoria").value;

if (filtro !== "" &&
    material.categoriaId !== Number(filtro)) {
  continue;
}
\end{lstlisting}

  \texttt{continue} pula para o próximo material\\
  e não cria o \texttt{li}.
\end{frame}

\begin{frame}{Desafio 2 --- não apagar categoria em uso}
  \onde{\texttt{CategoriaController.java} e \texttt{app.js}}

  Antes de remover a categoria, percorram
  \texttt{Dados.materiais}.

  \vspace{0.4em}
  Se algum material tiver aquele \texttt{categoriaId},\\
  lancem \texttt{ResponseStatusException} com
  \texttt{HttpStatus.BAD\_REQUEST}.

  \vspace{0.4em}
  No JS, o clique de excluir categoria usa \texttt{catch}\\
  e escreve em \texttt{\#mensagem}:\\
  ``Categoria em uso''.

  \vspace{0.4em}
  Testem apagar Papelaria com a Caneta ainda cadastrada:\\
  tem que recusar.
\end{frame}

\begin{frame}{Desafio 3 --- etiquetas (N{:}N)}
  \onde{Arquivos novos no laboratorio}

  \begin{enumerate}
    \item Classe \texttt{Etiqueta}: \texttt{id} e \texttt{nome}
    \item Classe \texttt{MaterialEtiqueta}:
          \texttt{materialId} e \texttt{etiquetaId}
    \item Listas dessas duas classes em \texttt{Dados}
    \item GET e POST de \texttt{/api/etiquetas}
    \item POST \texttt{/api/materiais/\{id\}/etiquetas/\{etiquetaId\}}
          para vincular
  \end{enumerate}

  \vspace{0.4em}
  Seed sugerido: Fragil e Uso interno.
\end{frame}

\begin{frame}[fragile]{Desafio 3 --- o POST de vincular}
  O método não cria material nem etiqueta.\\
  Só acrescenta uma linha na lista do meio.

\begin{lstlisting}[language=Java,basicstyle=\ttfamily\scriptsize]
@PostMapping("/api/materiais/{id}/etiquetas/{etiquetaId}")
public MaterialEtiqueta vincular(@PathVariable int id,
    @PathVariable int etiquetaId) {
  MaterialEtiqueta ligacao = new MaterialEtiqueta();
  ligacao.setMaterialId(id);
  ligacao.setEtiquetaId(etiquetaId);
  Dados.ligacoes.add(ligacao);
  return ligacao;
}
\end{lstlisting}

  Construtor vazio + gets e sets na classe
  \texttt{MaterialEtiqueta}.
\end{frame}

\begin{frame}{Desafio 3 --- mostrar na lista}
  No JS, peçam GET das ligações (um GET
  \texttt{/api/ligacoes} ajuda).

  \vspace{0.4em}
  Para cada material, percorram as ligações.\\
  Quando o \texttt{materialId} bater, achem o nome da etiqueta\\
  e colem no texto do \texttt{li}.

  \vspace{0.4em}
  Um material pode mostrar duas etiquetas.\\
  A mesma etiqueta pode aparecer em dois materiais.

  \vspace{0.4em}
  Extra: botão para desvincular, apagando aquela linha\\
  da lista do meio --- sem apagar o material.
\end{frame}

\begin{frame}{Extra}
  Se terminarem os três, ainda no \textbf{laboratorio}:

  \vspace{0.4em}
  \begin{itemize}
    \item Não vincular a mesma etiqueta duas vezes
          no mesmo material
    \item Recusar vincular se o material ou a etiqueta
          não existir
    \item Página simples só para cadastrar categoria
  \end{itemize}
\end{frame}

\begin{frame}{Checklist}
  \begin{itemize}
    \item [ ] Estou no projeto \texttt{laboratorio}?
    \item [ ] As listas estão em \texttt{Dados}?
    \item [ ] \texttt{/api/categorias} devolve as três?
    \item [ ] O material tem \texttt{categoriaId}?
    \item [ ] A lista mostra o \textbf{nome} da categoria?
    \item [ ] O PUT não zera a categoria?
  \end{itemize}
\end{frame}

\begin{frame}{Intervalo}
  \begin{center}
    \Large Retomamos às \textbf{20h20}\\[0.8em]
    \normalsize
    Fechem o \texttt{laboratorio}.\\
    Parem o Spring.\\
    Na segunda parte abrimos o \textbf{StockSales}.
  \end{center}
\end{frame}

% #############################################################################
\section{Parte 2 --- StockSales}

\begin{frame}{Mudamos de projeto}
  \onde{Agora é o projeto \textbf{StockSales}.\\
        O laboratorio já está encerrado.}

  Parem o Spring do laboratorio.\\
  Abram a pasta do StockSales e subam de novo.
\end{frame}

\begin{frame}{Onde o StockSales está}
  Até agora a loja tem \textbf{só} produto.

  \vspace{0.4em}
  \begin{itemize}
    \item Um model: \texttt{Produto}
    \item Um Controller: GET, POST, PUT, DELETE
    \item A lista fica na classe do Controller
    \item Tela de produtos com formulário
  \end{itemize}

  \vspace{0.5em}
  Ainda \textbf{não} tem cliente.\\
  Ainda \textbf{não} tem venda.\\
  Ainda \textbf{não} tem a classe \texttt{Dados}.
\end{frame}

\begin{frame}{O que vamos fazer no StockSales}
  Aplicar o 1{:}N da primeira parte na loja.

  \vspace{0.5em}
  \begin{enumerate}
    \item Classe \texttt{Dados} --- as listas num lugar só
    \item Cadastro de \textbf{cliente}
    \item Cadastro de \textbf{venda}: um cliente e um produto
  \end{enumerate}

  \vspace{0.5em}
  A venda guarda \texttt{clienteId} e \texttt{produtoId}.\\
  É o mesmo \texttt{categoriaId} do material, com outro nome.
\end{frame}

\begin{frame}{O desenho da venda de hoje}
  \begin{center}
  \begin{tikzpicture}[
      box/.style={draw, rounded corners, align=center,
                  minimum width=2.6cm, minimum height=0.9cm,
                  font=\small},
      seta/.style={-{Latex}, thick}
    ]
    \node[box, fill=blue!15] (cli) {Cliente};
    \node[box, fill=yellow!25, right=2.0cm of cli] (ven) {Venda};
    \node[box, fill=green!20, right=2.0cm of ven] (pro) {Produto};

    \draw[seta] (cli) -- node[above, font=\scriptsize]{1{:}N} (ven);
    \draw[seta] (pro) -- node[above, font=\scriptsize]{1{:}N} (ven);
  \end{tikzpicture}
  \end{center}

  \vspace{0.6em}
  Um cliente tem várias vendas.\\
  Um produto aparece em várias vendas.\\
  Cada venda, \textbf{nesta aula}, aponta para \textbf{um} produto.
\end{frame}

\begin{frame}[fragile]{Classe Dados no StockSales}
  \onde{Projeto \textbf{StockSales}\\
        Arquivo novo:\\
        \texttt{src/main/java/com/entra21/stocksales/Dados.java}}

\begin{lstlisting}[language=Java,basicstyle=\ttfamily\scriptsize]
package com.entra21.stocksales;

import java.util.ArrayList;
import com.entra21.stocksales.model.Produto;

public class Dados {
  public static ArrayList<Produto> produtos = new ArrayList<>();
  public static int proximoIdProduto = 1;
}
\end{lstlisting}

  \texttt{static} = uma lista só.\\
  Cliente e venda entram nessa classe daqui a pouco.
\end{frame}

\begin{frame}[fragile]{Construtor cheio do Produto}
  \onde{\texttt{model/Produto.java}}

  Se ainda não tiver, acrescentem. Mantenham o vazio.

\begin{lstlisting}[language=Java,basicstyle=\ttfamily\footnotesize]
public Produto(int id, String nome, double preco, int estoque) {
  this.id = id;
  this.nome = nome;
  this.preco = preco;
  this.estoque = estoque;
}
\end{lstlisting}

  Sem o vazio, o POST quebra.\\
  Sem o cheio, o seed do construtor fica longo.
\end{frame}

\begin{frame}[fragile]{Mover a lista de produtos}
  \onde{\texttt{ProdutoController.java}}

  Tirem \texttt{produtos} e \texttt{proximoID} da classe.

\begin{lstlisting}[language=Java,basicstyle=\ttfamily\scriptsize]
public ProdutoController() {
  if (Dados.produtos.isEmpty()) {
    Dados.produtos.add(new Produto(1, "Teclado Mecanico", 250, 12));
    Dados.produtos.add(new Produto(2, "Mouse Gamer", 120, 30));
    Dados.produtos.add(new Produto(3, "Monitor 24", 900, 8));
    Dados.produtos.add(new Produto(4, "Headset", 200, 15));
    Dados.proximoIdProduto = 5;
  }
}
\end{lstlisting}

  Em GET, POST, PUT e DELETE: \texttt{Dados.produtos}\\
  e \texttt{Dados.proximoIdProduto}.
\end{frame}

\begin{frame}{Conferindo depois de mover}
  \begin{itemize}
    \item [ ] A página de produtos abre
    \item [ ] Teclado, Mouse, Monitor e Headset aparecem
    \item [ ] Cadastrar um produto novo atualiza a lista
  \end{itemize}

  \vspace{0.5em}
  O \texttt{if (Dados.produtos.isEmpty())} evita duplicar\\
  o seed se o Spring criar o Controller de novo.
\end{frame}

\begin{frame}[fragile,shrink=5]{Classe Cliente}
  \onde{Arquivo novo:\\
        \texttt{src/main/java/com/entra21/stocksales/model/Cliente.java}}

\begin{lstlisting}[language=Java,basicstyle=\ttfamily\scriptsize]
public class Cliente {
  private int id;
  private String nome;
  private String email;
  private String telefone;

  public Cliente() {}

  public Cliente(int id, String nome, String email, String telefone) {
    this.id = id;
    this.nome = nome;
    this.email = email;
    this.telefone = telefone;
  }
}
\end{lstlisting}

  Completem os gets e sets dos quatro atributos.
\end{frame}

\begin{frame}[fragile]{Clientes em Dados}
  \onde{\texttt{Dados.java}, acrescentar}

\begin{lstlisting}[language=Java,basicstyle=\ttfamily\footnotesize]
public static ArrayList<Cliente> clientes = new ArrayList<>();
public static int proximoIdCliente = 1;
\end{lstlisting}

  Import de \texttt{Cliente} no topo do arquivo.
\end{frame}

\begin{frame}[fragile,shrink=8]{ClienteController --- seed e GET}
  \onde{Arquivo novo:\\
        \texttt{controller/ClienteController.java}}

\begin{lstlisting}[language=Java,basicstyle=\ttfamily\scriptsize]
@RestController
public class ClienteController {

  public ClienteController() {
    if (Dados.clientes.isEmpty()) {
      Dados.clientes.add(new Cliente(1, "Ana", "ana@email.com", "1199"));
      Dados.clientes.add(new Cliente(2, "Bruno", "bruno@email.com", "1188"));
      Dados.clientes.add(new Cliente(3, "Carla", "carla@email.com", "1177"));
      Dados.proximoIdCliente = 4;
    }
  }

  @GetMapping("/api/clientes")
  public ArrayList<Cliente> listar() {
    return Dados.clientes;
  }
}
\end{lstlisting}
\end{frame}

\begin{frame}[fragile]{POST de cliente}
  \onde{Ainda em \texttt{ClienteController.java}}

\begin{lstlisting}[language=Java,basicstyle=\ttfamily\footnotesize]
@PostMapping("/api/clientes")
public Cliente cadastrar(@RequestBody Cliente cliente) {
  cliente.setId(Dados.proximoIdCliente);
  Dados.proximoIdCliente = Dados.proximoIdCliente + 1;
  Dados.clientes.add(cliente);
  return cliente;
}
\end{lstlisting}

  Testem \texttt{/api/clientes} no navegador:\\
  Ana, Bruno e Carla.
\end{frame}

\begin{frame}{Tela de clientes}
  \onde{Arquivos novos:\\
        \texttt{static/clientes.html}\\
        \texttt{static/js/clientes.js}}

  Copiem o desenho de produtos:

  \vspace{0.3em}
  \begin{itemize}
    \item Formulário: nome, e-mail, telefone
    \item Lista na \texttt{<ul>}
    \item \texttt{carregarClientes} com GET
    \item Submit com POST e \texttt{preventDefault}
  \end{itemize}

  \vspace{0.4em}
  Na home, um link para \texttt{/clientes.html}.
\end{frame}

\begin{frame}{O que é uma venda nesta aula}
  Uma venda tem \textbf{um} cliente e \textbf{um} produto.

  \vspace{0.5em}
  \begin{itemize}
    \item Quem compra: o cliente --- \texttt{clienteId}
    \item O quê: um produto --- \texttt{produtoId}
    \item Quantas unidades
    \item Total = preço $\times$ quantidade
  \end{itemize}

  \vspace{0.5em}
  Ao confirmar: o estoque daquele produto diminui.\\
  Se não houver estoque, a venda \textbf{não} é gravada.
\end{frame}

\begin{frame}[fragile]{Classe Venda}
  \onde{Arquivo novo:\\
        \texttt{src/main/java/com/entra21/stocksales/model/Venda.java}}

\begin{lstlisting}[language=Java,basicstyle=\ttfamily\footnotesize]
public class Venda {
  private int id;
  private int clienteId;
  private int produtoId;
  private int quantidade;
  private double total;

  public Venda() {
  }
}
\end{lstlisting}

  Construtor vazio + gets e sets dos cinco atributos.
\end{frame}

\begin{frame}[fragile]{Vendas em Dados}
  \onde{\texttt{Dados.java}}

\begin{lstlisting}[language=Java,basicstyle=\ttfamily\footnotesize]
public static ArrayList<Venda> vendas = new ArrayList<>();
public static int proximoIdVenda = 1;
\end{lstlisting}
\end{frame}

\begin{frame}[fragile]{O JSON da venda}
  O JavaScript manda só isto:

\begin{lstlisting}[basicstyle=\ttfamily\footnotesize]
{"clienteId": 1, "produtoId": 2, "quantidade": 3}
\end{lstlisting}

  \vspace{0.4em}
  O servidor faz o resto:
  \begin{enumerate}
    \item Achar o cliente
    \item Achar o produto
    \item Se \texttt{estoque < quantidade}, recusar
    \item \texttt{total = preco * quantidade}
    \item Diminuir o estoque
    \item Guardar a venda e devolver o JSON
  \end{enumerate}
\end{frame}

\begin{frame}[fragile]{VendaController --- achar cliente}
  \onde{Arquivo novo: \texttt{controller/VendaController.java}}

\begin{lstlisting}[language=Java,basicstyle=\ttfamily\scriptsize]
@PostMapping("/api/vendas")
public Venda cadastrar(@RequestBody Venda pedido) {
  Cliente cliente = null;
  for (int i = 0; i < Dados.clientes.size(); i++) {
    if (Dados.clientes.get(i).getId() == pedido.getClienteId()) {
      cliente = Dados.clientes.get(i);
    }
  }
  if (cliente == null) {
    throw new ResponseStatusException(
        HttpStatus.BAD_REQUEST, "Cliente nao encontrado");
  }
\end{lstlisting}

  O método ainda não acabou.
\end{frame}

\begin{frame}[fragile]{VendaController --- achar produto}
  \onde{Ainda no mesmo método \texttt{cadastrar}}

\begin{lstlisting}[language=Java,basicstyle=\ttfamily\scriptsize]
  Produto produto = null;
  for (int i = 0; i < Dados.produtos.size(); i++) {
    if (Dados.produtos.get(i).getId() == pedido.getProdutoId()) {
      produto = Dados.produtos.get(i);
    }
  }
  if (produto == null) {
    throw new ResponseStatusException(
        HttpStatus.BAD_REQUEST, "Produto nao encontrado");
  }
\end{lstlisting}
\end{frame}

\begin{frame}[fragile]{VendaController --- estoque e gravação}
  \onde{Ainda no mesmo método \texttt{cadastrar}}

\begin{lstlisting}[language=Java,basicstyle=\ttfamily\scriptsize]
  if (produto.getEstoque() < pedido.getQuantidade()) {
    throw new ResponseStatusException(
        HttpStatus.BAD_REQUEST, "Estoque insuficiente");
  }

  pedido.setId(Dados.proximoIdVenda);
  Dados.proximoIdVenda = Dados.proximoIdVenda + 1;
  pedido.setTotal(produto.getPreco() * pedido.getQuantidade());
  produto.setEstoque(
      produto.getEstoque() - pedido.getQuantidade());
  Dados.vendas.add(pedido);
  return pedido;
}
\end{lstlisting}
\end{frame}

\begin{frame}[fragile]{GET das vendas}
  \onde{Ainda em \texttt{VendaController.java}}

\begin{lstlisting}[language=Java,basicstyle=\ttfamily\footnotesize]
@GetMapping("/api/vendas")
public ArrayList<Venda> listar() {
  return Dados.vendas;
}
\end{lstlisting}

  \vspace{0.4em}
  Imports: \texttt{HttpStatus}, \texttt{ResponseStatusException},
  \texttt{Dados}, \texttt{Cliente}, \texttt{Produto}, \texttt{Venda}.
\end{frame}

\begin{frame}[fragile]{A tela usa o select da primeira parte}
  Em vez de a pessoa lembrar o id, a página mostra o nome.

\begin{lstlisting}[basicstyle=\ttfamily\footnotesize]
<select id="clienteId"></select>
<select id="produtoId"></select>
\end{lstlisting}

  \vspace{0.4em}
  O JS preenche as opções com GET de clientes\\
  e GET de produtos --- o mesmo \texttt{<option>}\\
  que vocês fizeram para categoria.
\end{frame}

\begin{frame}[fragile]{vendas-nova.html}
  \onde{\texttt{static/vendas-nova.html}, arquivo novo}

\begin{lstlisting}[basicstyle=\ttfamily\scriptsize]
<h1>Nova venda</h1>
<a href="/">Voltar</a>
<label>Cliente</label>
<select id="clienteId"></select>
<label>Produto</label>
<select id="produtoId"></select>
<label>Quantidade</label>
<input id="quantidade" type="number">
<button id="btn-vender" type="button">Confirmar venda</button>
<p id="mensagem"></p>
\end{lstlisting}

  Não esqueçam o Axios e \texttt{/js/vendas-nova.js}.
\end{frame}

\begin{frame}[fragile]{Preencher o select de clientes}
  \onde{\texttt{static/js/vendas-nova.js}, arquivo novo}

\begin{lstlisting}[basicstyle=\ttfamily\scriptsize]
function carregarClientes() {
  axios.get("/api/clientes").then(function (resposta) {
    var select = document.getElementById("clienteId");
    select.innerHTML = "";
    var clientes = resposta.data;
    for (var i = 0; i < clientes.length; i++) {
      var opcao = document.createElement("option");
      opcao.value = clientes[i].id;
      opcao.textContent = clientes[i].nome;
      select.appendChild(opcao);
    }
  });
}
\end{lstlisting}
\end{frame}

\begin{frame}[fragile]{Select de produtos}
  \onde{Ainda em \texttt{vendas-nova.js}}

\begin{lstlisting}[basicstyle=\ttfamily\scriptsize]
function carregarProdutos() {
  axios.get("/api/produtos").then(function (resposta) {
    var select = document.getElementById("produtoId");
    select.innerHTML = "";
    var produtos = resposta.data;
    for (var i = 0; i < produtos.length; i++) {
      var opcao = document.createElement("option");
      opcao.value = produtos[i].id;
      opcao.textContent = produtos[i].nome
        + " --- estoque: " + produtos[i].estoque;
      select.appendChild(opcao);
    }
  });
}
\end{lstlisting}
\end{frame}

\begin{frame}[fragile]{Confirmar a venda --- o JSON}
  \onde{Ainda em \texttt{vendas-nova.js}}

\begin{lstlisting}[basicstyle=\ttfamily\scriptsize]
document.getElementById("btn-vender")
  .addEventListener("click", function () {
    var dados = {
      clienteId: Number(document.getElementById("clienteId").value),
      produtoId: Number(document.getElementById("produtoId").value),
      quantidade: Number(document.getElementById("quantidade").value)
    };
\end{lstlisting}

  Os três valores vão como número, não como texto.
\end{frame}

\begin{frame}[fragile]{Confirmar a venda --- enviar}
  \onde{Ainda no clique de Confirmar}

\begin{lstlisting}[basicstyle=\ttfamily\scriptsize]
    axios.post("/api/vendas", dados)
      .then(function () {
        document.getElementById("mensagem").textContent =
          "Venda registrada";
        carregarProdutos();
      })
      .catch(function () {
        document.getElementById("mensagem").textContent =
          "Nao foi possivel registrar a venda";
      });
  });

carregarClientes();
carregarProdutos();
\end{lstlisting}
\end{frame}

\begin{frame}[fragile]{Histórico de vendas}
  \onde{\texttt{static/vendas.html} e \texttt{static/js/vendas.js}}

\begin{lstlisting}[basicstyle=\ttfamily\scriptsize]
axios.get("/api/vendas").then(function (resposta) {
  var vendas = resposta.data;
  var lista = document.getElementById("lista-vendas");
  lista.innerHTML = "";
  for (var i = 0; i < vendas.length; i++) {
    var li = document.createElement("li");
    li.textContent = "Cliente " + vendas[i].clienteId
      + " - produto " + vendas[i].produtoId
      + " - qtd " + vendas[i].quantidade
      + " - total R$ " + vendas[i].total;
    lista.appendChild(li);
  }
});
\end{lstlisting}

  Na home: links para clientes, nova venda e histórico.
\end{frame}

\begin{frame}{O caminho da primeira venda}
  \begin{enumerate}
    \item A pessoa escolhe Ana, Mouse Gamer, quantidade 2
    \item JS manda POST \texttt{/api/vendas}
    \item Servidor acha cliente e produto em \texttt{Dados}
    \item Confere estoque, calcula total, diminui estoque
    \item A tela de produtos, ao recarregar, já mostra
          o estoque menor
  \end{enumerate}
\end{frame}

\begin{frame}{Conferindo no StockSales}
  \begin{itemize}
    \item [ ] \texttt{/api/clientes} lista Ana, Bruno e Carla
    \item [ ] Dá para cadastrar um cliente novo
    \item [ ] Nova venda baixa o estoque
    \item [ ] Vender acima do estoque mostra a mensagem
    \item [ ] Histórico lista a venda com o total
    \item [ ] Home leva a produtos, clientes, nova venda
          e histórico
  \end{itemize}
\end{frame}

\begin{frame}{Extra no StockSales}
  Se der tempo --- é o N{:}N da primeira parte:

  \vspace{0.4em}
  \begin{itemize}
    \item Vários produtos na mesma venda, com uma lista
          \texttt{ItemVenda} (\texttt{produtoId} + \texttt{quantidade})
    \item No histórico, mostrar o \textbf{nome} do cliente
          e do produto, não só o id
    \item PUT e DELETE de cliente
    \item Não permitir quantidade 0 ou negativa
  \end{itemize}
\end{frame}

\begin{frame}{O que não entra hoje}
  \begin{itemize}
    \item Cancelar venda e devolver estoque
    \item Banco de dados
    \item Login
    \item Anotações JPA
  \end{itemize}

  \vspace{0.5em}
  \begin{block}{Meta no StockSales}
    Cliente cadastrável.\\
    Uma venda de um item, com estoque baixando.
  \end{block}
\end{frame}

\begin{frame}{Critérios de aceite}
  \begin{itemize}
    \item [ ] Cliente aparece na lista e no select da venda
    \item [ ] Venda grava \texttt{clienteId} e \texttt{produtoId}
    \item [ ] Total é calculado no servidor
    \item [ ] Estoque diminui
    \item [ ] Estoque insuficiente recusa a venda
  \end{itemize}
\end{frame}

% #############################################################################
\section{Recapitulando}

\begin{frame}{O que vimos hoje}
  \begin{enumerate}
    \item Relacionamento é um cadastro apontando para outro
    \item 1{:}N: o lado ``muitos'' guarda o id do lado ``um''
    \item N{:}N: uma terceira lista guarda os dois ids
    \item O \texttt{select} mostra o nome e manda o id
    \item No StockSales, a venda aponta para cliente e produto
  \end{enumerate}
\end{frame}

\begin{frame}{Os dois desenhos, lado a lado}
  \begin{center}
  \begin{tabular}{ll}
    \toprule
    laboratorio & StockSales \\
    \midrule
    Categoria 1{:}N Material & Cliente 1{:}N Venda \\
    \texttt{categoriaId} no material & \texttt{clienteId} na venda \\
    Material aponta para categoria & Venda aponta para produto \\
    Etiqueta N{:}N Material & extra: vários itens na venda \\
    \bottomrule
  \end{tabular}
  \end{center}

  \vspace{0.6em}
  Mesma ideia. Nomes diferentes.
\end{frame}

\begin{frame}{Na próxima aula}
  \begin{block}{Aula 06}
    \begin{itemize}
      \item Validação dos dados que chegam
      \item Erros HTTP com mensagem mais clara
      \item Organizar o que já cresceu no StockSales
    \end{itemize}
  \end{block}

  \vspace{0.6em}
  Tragam o StockSales com venda de um item funcionando.
\end{frame}

\begin{frame}
  \begin{center}
    \Huge Obrigado!\\[1.2em]
    \normalsize Mauricio Duailibi Neto\\
    Turma Java Entra21
  \end{center}
\end{frame}

\end{document}
